\section*{Chapitre \ref{ch:b3:climate-models} --- Modèles climatiques}

\begin{solution}{exo:b3:climate-models:1}
$A = 14.0 - 1.5 = \num{12.5}$ et $B = 0.14 - 0.05 = \num{0.09}$, en
\unit{kcal.cm^{-2}} par mois. Un mois vaut \qty{2.63e6}{s}, donc
\qty{1}{kcal.cm^{-2}} par mois vaut $\num{4.184e7}/\num{2.63e6} =
\qty{15.91}{W/m^2}$~: $A = \qty{199}{W/m^2}$, $B =
\qty{1.43}{W.m^{-2}.\celsius^{-1}}$. À \qty{15}{\celsius}, $199 + 21.5 =
\qty{220}{W/m^2}$.
\end{solution}

\begin{solution}{exo:b3:climate-models:2}
$(238.0 - 211.9)/1.816 = \qty{14.4}{\celsius}$~; avec 0.31, $(234.6 -
211.9)/1.816 = \qty{12.5}{\celsius}$. Chaque 0.01 d'albédo coûte
\qty{3.4}{W/m^2}, et $3.4/1.816 = \qty{1.9}{\celsius}$.
\end{solution}

\begin{solution}{exo:b3:climate-models:3}
{\sloppy $C = 1025\times3992\times60 = \qty{2.46e8}{J.m^{-2}.K^{-1}}$~; divisé par
les \qty{3.156e7}{s} d'une année, \qty{7.78}{W.an.m^{-2}.K^{-1}}. Alors $\tau
= 7.78/1.816 = \num{4.3}~\text{ans}$, et $7.78/1.33 = \num{5.8}~\text{ans}$.\par}
\end{solution}

\begin{solution}{exo:b3:climate-models:4}
ECS $= 3.93/1.33 = \qty{2.95}{\celsius}$~; TCR $\approx 3.93/1.95 =
\qty{2.02}{\celsius}$~; rapport \num{0.68}.
\end{solution}

\begin{solution}{exo:b3:climate-models:5}
$Q = \qty{335}{W/m^2}$. Boule de neige~: $(335\times0.38 - 211.9)/1.816 =
\qty{-46.6}{\celsius}$~; chaud~: $(335\times0.70 - 211.9)/1.816 =
\qty{12.4}{\celsius}$, encore au-dessus de \qty{10}{\celsius}~; rampe~: $335(0.54 +
0.016\,T) = 211.9 + 1.816\,T$ donne $T = \qty{8.7}{\celsius}$. $N' =
\qty{-1.816}{W.m^{-2}.\celsius^{-1}}$ aux deux premiers, stable, et $-1.816
+ 335\times0.016 = \num{+3.54}$ à la rampe, instable.
\end{solution}

\begin{solution}{exo:b3:climate-models:6}
La branche chaude se termine là où sa température atteint \qty{0}{\celsius}~:
$0.70\,Q = A$, $Q = \qty{302.7}{W/m^2}$, soit une constante solaire de
\qty{1211}{W/m^2}, \qty{11.0}{\percent} en dessous de celle d'aujourd'hui. La boule de neige fond
là où $0.38\,Q = A - 20B$, $Q = \qty{462.1}{W/m^2}$, $S_0 =
\qty{1848}{W/m^2}$. Le seuil vaut $B(T^* - T_F)/(Q_0(1 - \alpha_F))$,
soit l'écart entre la température actuelle $T^*$ et la température à laquelle la
glace commence à s'étendre, ramené à l'ensoleillement~: $1.816\times14.37/238 =
\qty{11.0}{\percent}$ ici, $1.816\times4.37/238 = \qty{3.3}{\percent}$ dans
le texte.
\end{solution}

\begin{solution}{exo:b3:climate-models:7}
(a)~\qty{-45.5}{\celsius}. (b)~$340\times0.38 + F = 211.9 - 18.2$~: $F =
\qty{64.5}{W/m^2}$, soit 16 fois le forçage d'un doublement. (c)~Tant que le
forçage dure, la planète bascule sur la branche chaude, $(238.0 + 64.5 -
211.9)/1.816 = \qty{50}{\celsius}$~; une fois celui-ci supprimé, la planète se refroidit
le long de la branche chaude jusqu'à \qty{14.4}{\celsius}, qui existe avec le Soleil d'aujourd'hui~:
elle ne retourne pas à la boule de neige.
\end{solution}

\begin{solution}{exo:b3:climate-models:8}
$\Sigma = 1.80/7.0 + 0.70/90 = \qty{0.2649}{an^{-1}}$, $\Pi = 1.10\times0.70/
(7.0\times90) = \qty{1.222e-3}{an^{-2}}$. Les racines donnent $\tau_f =
\num{3.84}~\text{ans}$ et $\tau_s = \num{213}~\text{ans}$~; les approximations,
$7.0/1.80 = \num{3.89}\text{ ans}$ et $90\times1.80/(1.10\times0.70) =
\num{210}~\text{ans}$. $a_f = 0.597$ exactement, $1.10/1.80 = 0.611$
environ. ECS $= 3.93/1.10 = \qty{3.57}{\celsius}$~; TCR $\approx
3.93/1.80 = \qty{2.18}{\celsius}$.
\end{solution}

\begin{solution}{exo:b3:climate-models:9}
Dénominateur $B + \beta = \num{5.555}$. Équateur~: $(340\times1.221\times0.70 -
211.9 + 3.74\times12.81)/5.555 = \qty{22.8}{\celsius}$~; pôle~:
$(340\times0.506\times0.38 - 211.9 + 47.9)/5.555 = \qty{-17.7}{\celsius}$.
Le transport retire $3.74\times(22.8 - 12.8) = \qty{37}{W/m^2}$ à
l'équateur et apporte $3.74\times30.6 = \qty{114}{W/m^2}$ au pôle.
\end{solution}

\begin{solution}{exo:b3:climate-models:10}
(a)~$(1 - x^2)P_2' = 3x - 3x^3$, dont la dérivée vaut $3 - 9x^2 = -6P_2$.
(b)~Le terme de transport vaut $-\mathcal D T_2(-6P_2) = 6\mathcal DT_2P_2$.
Le bilan $Q(1 + s_2P_2)(1 - \alpha) - A - B(T_0 + T_2P_2) - 6\mathcal
DT_2P_2 = 0$ est vérifié pour tout $x$ si sa partie constante et sa partie en $P_2$
s'annulent, ce qui donne $T_0$ et $T_2$. (c)~$T(0) - T(1) = T_2(-\tfrac12 - 1)
= -1.5\,T_2$. Sans transport, $T_2 = 340\times(-0.476)\times0.70/1.816 =
\qty{-62.4}{\celsius}$~: un écart de \qty{94}{\celsius}. (d)~$T_2 =
-55.3/1.5 = \qty{-36.9}{\celsius}$, donc $B + 6\mathcal D = 113.3/36.9 =
3.07$ et $\mathcal D = \qty{0.21}{W.m^{-2}.\celsius^{-1}}$.
\end{solution}

\begin{solution}{exo:b3:climate-models:11}
(a)~$r = \tfrac12\Sigma\big(1 \pm \sqrt{1 - 4\Pi/\Sigma^2}\big)$~; avec
$\sqrt{1 - \epsilon} \approx 1 - \epsilon/2$, les racines valent $\Sigma -
\Pi/\Sigma \approx \Sigma$ et $\Pi/\Sigma$. (b)~Avec $T_d = 0$, $C\,\mathrm
dT/\mathrm dt = kt - (\gamma - \lambda)T$, dont la solution à partir du repos est $T =
k(t - \tau_f)/(\gamma - \lambda) + k\tau_f\mathrm e^{-t/\tau_f}/(\gamma -
\lambda)$, $\tau_f = C/(\gamma - \lambda)$. Pour $t \gg \tau_f$, $T \approx
kt/(\gamma - \lambda)$, et au doublement, $kt = F_{2\times}$~: TCR
$\approx F_{2\times}/(\gamma - \lambda)$. (c)~$\Sigma = 1.95/8.1 +
0.62/110 = \qty{0.2464}{an^{-1}}$, $\Pi = 1.33\times0.62/(8.1\times110) =
\qty{9.25e-4}{an^{-2}}$~; $\tau_f = \num{4.12}~\text{ans}$ et $\tau_s =
\num{262}~\text{ans}$ (environ $1/\Sigma = \num{4.06}~\text{ans}$ et
$\Sigma/\Pi = \num{266}\text{ ans}$)~; $4\Pi/\Sigma^2 = 0.061$.
\end{solution}

\begin{solution}{exo:b3:climate-models:12}
(a)~Moyenne \qty{3.50}{\celsius}, écart type \qty{0.81}{\celsius}.
(b)~ECS $= 1.88 + 1.86\,X$~; corrélation \num{0.992}~; écart type résiduel
\qty{0.11}{\celsius}. (c)~$1.88 + 1.86\times0.7 =
\qty{3.18}{\celsius}$, avec une dispersion $\sqrt{(1.86\times0.1)^2 + 0.11^2} =
\qty{0.22}{\celsius}$, soit environ un quart de la dispersion de l'ensemble.
(d)~Aucun mécanisme physique ne relie $X$ aux rétroactions~; un $X$ trouvé en
essayant de nombreuses observables, de sorte que la corrélation peut être due au hasard~; un biais
partagé par les huit modèles~; des modèles de la même famille comptés comme
indépendants.
\end{solution}

\begin{solution}{pb:b3:climate-models:1}
\textbf{1.} $Q_0 = 1360/4 = \qty{340}{W/m^2}$~: la planète intercepte le
rayonnement solaire sur un disque d'aire $\pi R^2$ et le répartit sur une sphère d'aire
$4\pi R^2$.

\textbf{2.} $340\times0.70 = \qty{238.0}{W/m^2}$~; $211.9 + 1.816\times15 =
\qty{239.1}{W/m^2}$.

\textbf{3.} $(238.0 - 211.9)/1.816 = \qty{14.4}{\celsius}$.

\textbf{4.} $(129.2 - 211.9)/1.816 = \qty{-45.5}{\celsius}$, en dessous de
\qty{-10}{\celsius}.

\textbf{5.} $340(0.54 + 0.016\,T) = 211.9 + 1.816\,T$~: $T =
(183.6 - 211.9)/(1.816 - 5.44) = \qty{7.8}{\celsius}$.

\textbf{6.} $(238.0/\num{5.670e-8})^{1/4} = \qty{254.5}{K}$,
\qty{-18.6}{\celsius}~: \qty{33}{\celsius} sous l'équilibre chaud, soit
l'effet de serre de la planète, intégré dans $A$ et $B$.

\textbf{7.} $N' = -B = \qty{-1.816}{W.m^{-2}.\celsius^{-1}}$ pour la
Terre boule de neige et l'équilibre chaud, tous deux stables~; $-1.816 + 5.44 =
\num{+3.62}$ pour l'équilibre intermédiaire, instable.

\textbf{8.} $\tau = 8.1/3.62 = \num{2.2}\text{ ans}$.

\textbf{9.} $8.1/1.816 = \num{4.5}\text{ ans}$.

\textbf{10.} $3.93/1.816 = \qty{2.2}{\celsius}$, en dessous de la meilleure
estimation évaluée de \qty{3}{\celsius} et en dessous de sa fourchette \emph{probable}, de 2.5 à
\qty{4}{\celsius}~: le $B$ du modèle, ajusté sur les latitudes, n'est pas
le paramètre de rétroaction climatique de l'ensemble du climat.

\textbf{11.} $\mathrm e^{-t/\tau} = 0.05$ pour $t = 3\tau =
\num{13}~\text{ans}$. Le climat réel réchauffe aussi l'océan profond, dont le
temps caractéristique se mesure en siècles.

\textbf{12.} Pour $T_F = \qty{10}{\celsius}$~; $Q = (211.9 + 18.2)/0.70 =
\qty{328.7}{W/m^2}$.

\textbf{13.} $S_0 = \qty{1315}{W/m^2}$, \qty{3.3}{\percent} sous celle d'aujourd'hui.

\textbf{14.} Vers la Terre boule de neige, $(328.7\times0.38 - 211.9)/1.816 =
\qty{-47.9}{\celsius}$.

\textbf{15.} $Q = (211.9 - 18.2)/0.38 = \qty{509.8}{W/m^2}$, $S_0 =
\qty{2039}{W/m^2}$, \qty{50}{\percent} au-dessus de celle d'aujourd'hui.

\textbf{16.} En descendant la branche chaude jusqu'à \qty{10}{\celsius} pour
\qty{1315}{W/m^2}, une chute à \qty{-48}{\celsius}, puis remontée par la
branche de la Terre boule de neige jusqu'à \qty{-45.5}{\celsius} pour \qty{1360}{W/m^2}~: la planète
reste gelée. Ce n'est qu'au-delà de \qty{2039}{W/m^2} qu'elle repartirait vers la
branche chaude.

\textbf{17.} $\bar\alpha = 0.30\times0.9738 + 0.62\times0.0262 =
\num{0.3084}$~; $\bar T = (340\times0.6916 - 211.9)/1.816 =
\qty{12.80}{\celsius}$.

\textbf{18.} $T_c = (340\times0.5643\times0.50 - 211.9 + 3.739\times12.80)/
5.555 = \qty{-12.26}{\celsius}$, \qty{3.5}{\celsius} plus chaud que la
réanalyse à \ang{72}.

\textbf{19.} $5.555\times(-12.26) + 211.9\times3.059 = \num{580.1}$, avec
$\beta/B = \num{2.059}$.

\textbf{20.} Les dénominateurs valent 1.7053, 1.7473, 1.7565, 1.7488 et
1.7362, et $Q/Q_0 = \num{1.0005}$, \num{0.9764}, \num{0.9713},
\num{0.9756} et \num{0.9827}.

\textbf{21.} Près de \ang{47.5}, à \num{0.971}. Les lignes de glace situées du
côté des pôles par rapport au minimum sont stables, $\mathrm dQ/\mathrm dx_s > 0$~; celles
situées vers l'équateur sont instables.

\textbf{22.} $Q/Q_0 = 1$ également entre \ang{30} et \ang{40}, proche de
\ang{30} (\ang{30.2} par interpolation)~: une ligne de glace instable. À \ang{0},
$Q/Q_0 = 580.1/(0.6106 + 0.7824)/340 = \num{1.225}$~: la Terre boule de neige tient
pour tout Soleil jusqu'à \qty{22}{\percent} plus brillant que le Soleil actuel, elle
est donc possible aujourd'hui.

\textbf{23.} \qty{2.9}{\percent}, contre \qty{3.3}{\percent} pour le
modèle zéro-dimensionnel et \qty{1.6}{\percent} pour Budyko~; la glace se situe
près de \ang{48} au seuil, contre \ang{50} chez Budyko. Le modèle latitudinal
se couvre de glace plus tôt que le modèle global, car le rayonnement solaire est le plus faible là où
la glace se forme.

\textbf{24.} La planète du modèle se couvre de glace pour une baisse de
\textbf{\qty[text-series-to-math=true]{2.9}{\percent}} de la
constante solaire (\qty{2.87}{\percent} d'après le calcul).
\end{solution}
